\unnumberedchapter{MỞ ĐẦU}

\renewcommand{\thetable}{MĐ.\arabic{table}}\renewcommand{\thefigure}{MĐ.\arabic{figure}}

\section*{1. Lý do chọn đề tài}

Trí tuệ nhân tạo (AI) đã đi một chặng đường dài từ những chương trình chứng minh định lý của thập niên 1950 tới các hệ thống tác tử dựa trên mô hình ngôn ngữ lớn hiện nay. Mỗi giai đoạn để lại một lớp kỹ thuật vẫn đang được dùng: luật chuyên gia, học máy thống kê, mạng nơ-ron tích chập và hồi quy, đồ thị tri thức, học sâu trên đồ thị và hệ đa tác tử. Một hệ thống thông minh thực tế hiếm khi chỉ dùng một kỹ thuật duy nhất. Nó thường là sự phối hợp của nhiều mô hình chuyên biệt, mỗi mô hình mạnh ở một loại dữ liệu, được điều phối để cùng đưa ra một quyết định có căn cứ.

Đái tháo đường (ĐTĐ) là một bài toán điển hình cho cách tiếp cận đó. Theo IDF Diabetes Atlas \cite{idf2021}, năm 2021 có khoảng 537 triệu người trưởng thành mắc ĐTĐ, và gần một nửa chưa được chẩn đoán. Dữ liệu để sàng lọc lại rất đa dạng: chỉ số tự khai (tuổi, BMI, tiền sử gia đình) phù hợp với học máy trên bảng; ảnh đáy mắt dùng để phát hiện biến chứng võng mạc là đất của CNN; chuỗi đo đường huyết liên tục cần mô hình tuần tự như RNN; hướng dẫn lâm sàng (ADA 2024 \cite{ada2024}, Quyết định 5481/QĐ-BYT \cite{byt5481}) là tri thức văn bản cần RAG và đồ thị tri thức; y văn có cấu trúc trích dẫn là dữ liệu đồ thị. Không mô hình đơn lẻ nào xử lý được tất cả. Cần một hệ đa tác tử phân công, đối chất và tổng hợp các kết luận.

Báo cáo này trình bày tiến trình phát triển của AI qua bảy chủ đề của môn học. Mỗi chủ đề được kiểm chứng bằng một thành phần chạy thật của cùng một hệ thống: \textbf{chatbot đa tác tử hỗ trợ sàng lọc nguy cơ đái tháo đường}, triển khai công khai tại \href{https://\deploydomain}{\texttt{\deploydomain}}. Cách làm này giúp lý thuyết của từng chương có một chỗ đứng cụ thể trong sản phẩm, và các chương nối với nhau như các bộ phận của một hệ thống.

\section*{2. Mục tiêu}

\begin{enumerate}
  \item Hệ thống hóa lịch sử AI và các kỹ thuật nền tảng: học máy cơ bản, CNN, RNN, đồ thị tri thức và RAG, mạng nơ-ron đồ thị, hệ đa tác tử. Với mỗi mô hình, giải thích chi tiết từng lớp (layer) và hàm hợp tạo nên nó.
  \item Với mỗi chương, chọn một bộ dữ liệu công khai liên quan tới ĐTĐ. Mô tả đường dẫn, kích thước, mẫu dữ liệu và phân bố, kèm nhận xét về chất lượng dữ liệu.
  \item Cài đặt và so sánh các mô hình trên cùng một giao thức đánh giá, kèm mô hình đối chứng. Báo cáo bằng thước đo phù hợp với từng bài toán (ROC-AUC, QWK, RMSE và lưới Clarke, Hit@k và MRR, độ chính xác phân loại nút).
  \item Đóng gói mô hình tốt nhất của mỗi chương thành dịch vụ web và triển khai qua Cloudflare Tunnel, để có thể kiểm tra trực tiếp từ Internet.
\end{enumerate}

\section*{3. Đối tượng, phạm vi và dữ liệu}

Đối tượng nghiên cứu là các mô hình và thuật toán của môn Hệ thống thông minh, áp dụng vào bài toán sàng lọc ĐTĐ. \tabref{tab:datasets} liệt kê bảy bộ dữ liệu được dùng. Tất cả đều công khai và tải được bằng mã nguồn kèm theo. Thực nghiệm được thiết kế để chạy trên một máy tính xách tay (CPU Intel i5-11400H, 16 GB RAM, GPU RTX 3050 4 GB). Mục tiêu là kiểm chứng lý thuyết, không phải đạt kết quả tốt nhất hiện có. Hệ thống chỉ phục vụ học tập và nghiên cứu, \emph{không thay thế chẩn đoán của bác sĩ}.

\begin{table}[H]
\centering
\caption{Bảy bộ dữ liệu dùng trong báo cáo}
\label{tab:datasets}
\small
\begin{tabularx}{\linewidth}{@{}L{1.4cm}L{3.6cm}L{3.4cm}Y@{}}
\toprule
\textbf{Chương} & \textbf{Bộ dữ liệu} & \textbf{Quy mô} & \textbf{Bài toán} \\
\midrule
1 & Pima Indians Diabetes \cite{smith1988} & 768 người, 8 đặc trưng & Phân loại ĐTĐ: luật chuyên gia → perceptron → MLP \\
2 & NHANES 2011–2018 \cite{nhanes} & 20.646 người, 13 đặc trưng & Phân loại, hồi quy HbA1c, gom cụm \\
3 & RetinaMNIST \cite{yang2023medmnist} & 1.600 ảnh đáy mắt & Phân độ bệnh võng mạc ĐTĐ (5 mức) \\
4 & ShanghaiT1DM/T2DM \cite{zhao2023shanghai} & 112 bệnh nhân, 125 bản ghi CGM & Dự báo glucose 15–60 phút \\
5 & KB ADA 2024 + MedQuAD \cite{benabacha2019} & 36 đoạn + 1.192 hỏi–đáp & Truy xuất và sinh câu trả lời có trích dẫn \\
6 & Pubmed-Diabetes \cite{namata2012} & 19.717 bài báo, 44.324 cạnh & Phân loại nút trên đồ thị trích dẫn \\
7 & NHANES (tập kiểm tra) & 400 ca, 400 tin nhắn & Đánh giá đầu–cuối hệ đa tác tử \\
\bottomrule
\end{tabularx}
\end{table}

\section*{4. Phương pháp nghiên cứu}

Báo cáo kết hợp bốn phương pháp. \emph{Tổng hợp tài liệu}: hệ thống hóa kiến thức từ giáo trình và các bài báo gốc (Perceptron, LSTM, ResNet, GCN, GAT, RAG\ldots). \emph{Lịch sử – logic}: trình bày mỗi kỹ thuật như lời giải cho hạn chế của kỹ thuật đi trước. \emph{Thực nghiệm có đối chứng}: mỗi chương có ít nhất một mô hình nền đơn giản (persistence, MLP không có tích chập, MLP không dùng đồ thị, chỉ BM25, một tác tử đơn lẻ\ldots). Dữ liệu được chia theo đơn vị độc lập (theo người, theo bệnh nhân) để tránh rò rỉ, siêu tham số chọn trên tập validation, còn tập kiểm tra chỉ dùng một lần. \emph{Kỹ nghệ hệ thống}: đóng gói mô hình thành API, kiểm thử tự động, và triển khai bằng Docker cùng Cloudflare Tunnel.

\section*{5. Kiến trúc tổng thể của hệ thống}

\figref{fig:overall} cho thấy vị trí của từng chương trong hệ thống. Người dùng trò chuyện với chatbot bằng văn bản tự do hoặc tải lên bệnh án. \emph{Orchestrator} (Chương 7) phân loại ý định rồi giao việc cho các tác tử chuyên biệt qua một \emph{MessageBus}. \emph{PredictionAgent} dùng mô hình học máy (Chương 2), CNN (Chương 3) và RNN (Chương 4). \emph{KnowledgeAgent} dùng RAG và đồ thị tri thức (Chương 5), cùng GNN để gán nhãn y văn (Chương 6). \emph{ConsensusCriticAgent} tổ chức tranh biện, liệt kê bằng chứng đồng thuận (+) và mâu thuẫn (−) trước khi đưa ra kết luận. Hệ chuyên gia luật ở Chương 1 tồn tại dưới dạng các ngưỡng ADA và lớp \emph{Safety Guardrail}. Đó là phần tri thức tường minh mà mọi mô hình học được đều phải tuân theo.

\begin{figure}[H]
\centering
\begin{tikzpicture}[font=\footnotesize,
  box/.style={draw, rounded corners=3pt, align=center, minimum height=10mm, text width=3.0cm, inner sep=3pt},
  ag/.style={box, fill=blue!6, draw=blue!50!black},
  ch/.style={font=\scriptsize\itshape, text=black!60},
  arr/.style={-{Stealth[length=2mm]}, thick, black!60}]
\node[box, fill=gray!10] (user) at (0,0) {Người dùng\\(chat, bệnh án, CGM)};
\node[ag] (orc) at (3.6,0) {Orchestrator\\ý định + kế hoạch};
\node[ag] (pred) at (7.6,1.5) {PredictionAgent\\ML · CNN · RNN};
\node[ag] (know) at (7.6,-1.5) {KnowledgeAgent\\RAG · KG · GNN};
\node[ag, fill=red!6, draw=red!50!black] (crit) at (11.6,0) {ConsensusCritic\\(+) / (−), Guardrail};
\node[box, fill=gray!10] (db) at (3.6,-2.9) {PostgreSQL\\pgvector · AGE};
\node[ch, below=1mm of orc] {Chương 7};
\node[ch, above=1mm of pred] {Chương 2, 3, 4};
\node[ch, below=1mm of know, xshift=9mm] {Chương 5, 6};
\node[ch, below=1mm of crit] {Chương 1, 7};
\draw[arr] (user) -- (orc);
\draw[arr] (orc.east) -- ++(0.5,0) |- (pred.west);
\draw[arr] (orc.east) -- ++(0.5,0) |- (know.west);
\draw[arr] (pred.east) -- ++(0.5,0) |- ([yshift=2mm]crit.west);
\draw[arr] (know.east) -- ++(0.5,0) |- ([yshift=-2mm]crit.west);
\draw[arr, dashed, red!60!black] (crit.north) |- node[pos=0.75, above, font=\scriptsize] {CHALLENGE} (pred.east);
\draw[arr, dashed, red!60!black] (crit.south) |- node[pos=0.75, below, font=\scriptsize] {CHALLENGE} (know.east);
\draw[arr] (db.east) -| ([xshift=-9mm]know.south);
\end{tikzpicture}
\caption{Vị trí các chương trong kiến trúc hệ thống chatbot đa tác tử}
\label{fig:overall}
\end{figure}

Toàn bộ hệ thống được đóng gói bằng Docker Compose với bốn dịch vụ: PostgreSQL 17 (có pgvector và Apache AGE), SeaweedFS (lưu file bệnh án), backend FastAPI (phục vụ cả API và giao diện) và \texttt{cloudflared}. \texttt{cloudflared} mở một kết nối đi ra (outbound) tới mạng Cloudflare, nhờ đó tên miền \texttt{\deploydomain} trỏ được về máy chủ mà không cần mở cổng hay có IP công khai \cite{cloudflaretunnel}. Mỗi chương có một endpoint riêng dưới \api{/api/lab/chX/...} và một thẻ trên trang thử nghiệm \api{/lab}. Chatbot đầy đủ nằm ở trang gốc \api{/}.

\section*{6. Bố cục của báo cáo}

Ngoài Mở đầu, Kết luận và Tài liệu tham khảo, báo cáo gồm bảy chương. \textbf{Chương 1} điểm lại lịch sử AI và tái hiện bốn thế hệ mô hình trên bộ Pima. \textbf{Chương 2} trình bày các thuật toán học máy cơ bản và so sánh tám bộ phân loại trên NHANES. \textbf{Chương 3} phân tích từng lớp của CNN và áp dụng vào phân độ bệnh võng mạc. \textbf{Chương 4} trình bày RNN, LSTM, GRU và bài toán dự báo đường huyết. \textbf{Chương 5} xây dựng đồ thị tri thức và đánh giá các chiến lược RAG. \textbf{Chương 6} trình bày GCN, GraphSAGE, GAT trên đồ thị trích dẫn y văn. \textbf{Chương 7} mô tả kiến trúc hệ đa tác tử, giao thức tranh biện và đánh giá đầu–cuối. Mỗi chương đều có bốn phần: cơ sở lý thuyết, dữ liệu, cài đặt và đánh giá, triển khai.
